% HD-F1-T37. GPT-6 Astra Pro, /root/transport; 22 November 2026 UTC.
\section[HD-F1-T37: odd symplectic level one]{HD-F1-T37: exclusion of every actual odd-characteristic symplectic level-one source, including all outer lifts}
\label{sec:f1-symplectic-odd-weil-source-exclusion}

In the original remaining T2-OTHER branch with almost simple
projective image, let the projective socle be
\(\operatorname{PSp}_{2r}(q)\), where \(q=p^f\) is odd and
\(r\ge5\). No irreducible constituent of the original coefficient
on its actual symplectic layer has level one. This statement
includes every original outer automorphism, scalar phase and
restriction multiplicity. In particular, the odd-characteristic
symplectic part of T36 has the sharpened exact range
\[
       2\le\ell\le3\cdot10^{30}-3
       \qquad(r\ge10^{31}).
\tag{F1T37.1}
\]
The proof gives the strict bound
\(\lvert\operatorname{Tr}_V(g)\rvert<3m/5\) for every
original nonscalar matrix in any proposed level-one case.

\paragraph{1. The actual layer and its automorphisms.}
Retain the original finite group \(G\subset GL_K(V)\), its
scalar subgroup \(Z\), the inverse image \(P\) of the projective
socle \(S\), and \(L=P'\). The retained actual maps give
\(P=ZL\), \(L=L'\), and \(L/(L\cap Z)=S\).
For \(r\ge5\), Malle--Testerman, Theorem~24.17,
Remark~24.19 and Tables~24.2--24.3 identify
\[
 H=\operatorname{Sp}_{2r}(q)\longrightarrow
 S=\operatorname{PSp}_{2r}(q)
\tag{F1T37.2}
\]
as the universal central cover, with centre \(\{I,-I\}\).
None of the listed multiplier or perfectness exceptions occurs.
Universality supplies the surjection \(\beta:H\to L\) used in
T16 and T36.

Steinberg's Theorem~30, with its complete finite-field proof
in steps~(1)--(4), gives the inner, diagonal, graph and field
automorphisms. The type \(C_r\) diagram has no nonidentity
graph automorphism in this range. More explicitly, each
automorphism induced by an original \(g\in G\) lifts to
\[
 \alpha(h)=a h^{p^e}a^{-1},\qquad
 a\in\operatorname{GSp}_{2r}(q),\quad 0\le e<f,
 \qquad \beta\alpha=c_g\beta.
\tag{F1T37.3}
\]
Here the field power is entrywise. The diagonal assertion in
this matrix description is literal: for simple roots
\(e_i-e_{i+1}\) and \(2e_r\), any prescribed root multipliers
\(d_1,\ldots,d_r\in\mathbf F_q^*\) are realized by
\[
 a=\operatorname{diag}(a_1,\ldots,a_r,
                  \lambda a_1^{-1},\ldots,\lambda a_r^{-1}),
 \quad a_r=1,\quad
 a_i=\prod_{j=i}^{r-1}d_j,\quad \lambda=d_r^{-1}.
\tag{F1T37.4}
\]
Indeed its root multipliers are \(a_i/a_{i+1}=d_i\) and
\(a_r^2/\lambda=d_r\); multiplying by an inner symplectic
matrix gives (F1T37.3). The two maps in its last equality
have the same projection to \(S\). Their pointwise quotient
is a homomorphism from the perfect group \(H\) to the central
kernel in \(L\), and hence is trivial. This also proves the
uniqueness of the lift used here.

An original nonscalar \(g\) induces a nonidentity automorphism
of \(S\), and therefore of \(H\). To check this assertion at
the actual projective image, its inclusion in
\(\operatorname{Aut}(S)\) has trivial centralizer of the inner
copy of \(S\): an automorphism commuting with every
\(c_s\) satisfies \(\alpha(s)s^{-1}\in Z(S)=1\).
Thus a matrix inducing the identity belongs to the original
scalar subgroup \(Z\).

\paragraph{2. Weyl operators and a genuine linear lift.}
Write \(A=\mathbf F_q^r\oplus\mathbf F_q^r\), with
\[
 \langle(x,y),(x',y')\rangle=x\cdot y'-x'\cdot y.
\]
For \(t\in\mathbf F_q^*\), put
\(\psi_t(z)=\exp(2\pi i\operatorname{Tr}_{q/p}(tz)/p)\)
and let \(\mathcal W_t=\mathbf C[\mathbf F_q^r]\), of
dimension \(D=q^r\). On functions define
\[
 (D_t(x,y)F)(u)=
       \psi_t(y\cdot u+x\cdot y/2)F(u+x).
\tag{F1T37.5}
\]
Direct multiplication and adjunction give
\[
 D_t(v)D_t(w)=\psi_t(\langle v,w\rangle/2)D_t(v+w),
 \quad D_t(v)^*=D_t(-v),
 \quad \operatorname{Tr}D_t(v)=D\,\mathbf1_{v=0}.
\tag{F1T37.6}
\]
For the trace assertion, a nonzero translation has no fixed
basis element; for zero translation the sum of a nontrivial
additive character is zero. Consequently the \(D^2\) operators
\(D_t(v)\) form an orthogonal basis of
\(\operatorname{End}(\mathcal W_t)\).

For \(h\in H\), sending \(D_t(v)\) to \(D_t(hv)\)
preserves multiplication, adjunction and the identity. It is
therefore a star automorphism of this full matrix algebra.
Every such automorphism is implemented by a unitary operator,
unique up to scalar: the images of the diagonal matrix units
are orthogonal rank-one projections; choose a unit vector in
the first image and transport it by the images of the other
matrix units to obtain the required orthonormal basis.
We obtain an actual homomorphism \(H\to PU(D)\).

Here is a direct linearization, so no uninspected cocycle
vanishing is needed. Pull back \(SU(D)\to PU(D)\), obtaining
a finite central extension \(E\to H\) with kernel
\(C\subset\mu_D\). Since \(H\) is perfect,
\(E=E'C\), and hence
\(E'=[E'C,E'C]=[E',E']\) is perfect and maps onto \(H\).
For \(z\in\ker(E'\to S)\), the commutator \([z,x]\)
lies in \(C\) for every \(x\in E'\). It is a homomorphism
\(E'\to C\), since \(C\) is central, and is trivial by
perfectness. Thus \(E'\to S\) is a central extension.
The universal property of (F1T37.2) gives \(H\to E'\)
over \(S\); its composition with \(E'\to H\) is the
identity, by uniqueness in the same universal property.
Projection to \(SU(D)\) gives a genuine unitary representation
\[
 R_t:H\longrightarrow U(\mathcal W_t),\qquad
 R_t(h)D_t(v)R_t(h)^{-1}=D_t(hv).
\tag{F1T37.7}
\]
Any two genuine lifts with this covariance differ by a
linear character of the perfect group \(H\), so agree.

\paragraph{3. The actual parity constituents and their labels.}
Let \((\Pi F)(u)=F(-u)\). Then
\(\Pi^2=1\), \(\operatorname{Tr}\Pi=1\), and
\(\Pi D_t(v)\Pi=D_t(-v)\). Conjugating \(\Pi\)
by any \(R_t(h)\) produces a scalar multiple of \(\Pi\);
its trace forces that scalar to be one. Thus its eigenspaces
\(\mathcal W_t^{\varepsilon}\), \(\varepsilon=\pm1\),
are actual \(H\)-submodules of dimensions
\[
              d_\varepsilon=(q^r+\varepsilon)/2.
\tag{F1T37.8}
\]
Conjugation on the basis (F1T37.6) is the permutation action
of \(H\) on \(A\). This action has exactly two orbits,
zero and nonzero: extend a nonzero vector to a symplectic
basis to obtain transitivity. Hence
\(\langle\operatorname{ch}R_t,\operatorname{ch}R_t\rangle_H=2\).
The two nonzero parity modules and the sum of squared
irreducible multiplicities show that they are irreducible
and inequivalent.

We verify the square-class labels as well. For a symmetric
matrix \(Q\), let \(u_Q(x,y)=(x,y-Qx)\) and
\[
          (M_QF)(u)=\psi_t(u^tQu/2)F(u).
\]
Its covariance is that of \(u_Q\), so
\(R_t(u_Q)=c(Q)M_Q\), where \(c\) is an additive
character on the symmetric matrices. The explicit Levi
operator \(F(u)\mapsto F(B^{-1}u)\), \(B\in GL_r(q)\),
shows that \(c\) is invariant under every congruence.
Using the nondegenerate trace pairing, write
\(c(Q)=\exp(2\pi i\operatorname{Tr}_{q/p}
\operatorname{tr}(CQ)/p)\), with \(C\) symmetric.
Invariance gives \(BCB^t=C\) for all \(B\).
For \(B=I\pm E_{ij}\), adding the two equations gives
\(C_{jj}=0\); subtracting and taking the \((i,i)\) entry
gives \(C_{ij}=0\). As \(i\ne j\) are arbitrary and
\(p\ne2\), \(C=0\). Therefore \(c=1\).

Take \(Q=\operatorname{diag}(1,0,\ldots,0)\).
With \(G_t=\sum_{x\in\mathbf F_q}\psi_t(x^2/2)\),
the trace on the two parity spaces is exactly
\[
 \operatorname{Tr}_{\mathcal W_t^\varepsilon}R_t(u_Q)
          =\frac{q^{r-1}G_t+\varepsilon}{2}.
\tag{F1T37.9}
\]
Indeed \(\operatorname{Tr}(\Pi M_Q)=1\). If \(\chi\)
is the quadratic character, counting the roots of \(x^2=z\)
gives \(G_t=\chi(t)G_1\). Moreover
\(\lvert G_t\rvert^2=q\): in the double sum for this
square, replace \((x,y)\) by \((x-y,x+y)\) and sum the
additive character of their product. Thus nonsquare ratios
of \(t\)'s give different characters in each parity.
For a square ratio \(t/t'=c^2\), the map
\(D_t(v)\mapsto D_{t'}(cv)\) is a star isomorphism.
Its implementing unitary intertwines the two genuine
\(H\)-actions by (F1T37.7) and perfectness; it also
intertwines \(\Pi\), since both parity operators have
trace one. These constructions give exactly four distinct
irreducibles, indexed by square class and parity.

These are the literal level-one constituents in GLT,
Definition~3.2. To specify conventions, the Heisenberg group
used in Guralnick--Magaard--Tiep, pp.~444--445, has product
\((z,v)(z',w)=(z+z'+\langle v,w\rangle,v+w)\).
The map \((z,v)\mapsto(z/2,v)\) identifies it with the
group for (F1T37.6), and changes its central parameter to
\(t/2\). This permutes the two square-class labels and
identifies the full sum \(\omega_r+\omega_r^*\) exactly.
Uniqueness of the genuine lift gives the module
identification, rather than just an equality of dimensions.

\paragraph{4. Every semilinear action on the four constituents.}
Let \(a\) and \(e\) be as in (F1T37.3), with multiplier
\(\lambda\), and put \(S_0=a\sigma_e\) on \(A\).
For a given \(t\), set
\[
                       t'=t^{p^e}/\lambda.
\tag{F1T37.10}
\]
Then
\[
 \operatorname{Tr}_{q/p}
       (t'\langle S_0v,S_0w\rangle)
 =\operatorname{Tr}_{q/p}
       (t^{p^e}\langle v,w\rangle^{p^e})
 =\operatorname{Tr}_{q/p}(t\langle v,w\rangle).
\tag{F1T37.11}
\]
Consequently \(D_t(v)\mapsto D_{t'}(S_0v)\) is an
actual star isomorphism. Its implementing unitary \(U\)
satisfies
\[
        U R_t(h)U^{-1}=R_{t'}(\alpha(h)),
        \qquad U\Pi U^{-1}=\Pi.
\tag{F1T37.12}
\]
For the first equality both sides are genuine
representations with the same covariance, so their scalar
quotient is a trivial character of \(H\). For the second,
covariance gives a scalar multiple and trace one fixes it.
Thus every original outer action preserves parity and
changes the square class precisely by that of \(\lambda\).
The Frobenius power preserves square classes. By
(F1T37.9), a constituent is fixed exactly when
\(\lambda\) is a square; a nonsquare exchanges the two
inequivalent constituents of its parity.

\paragraph{5. A trace-preserving model for each stabilizing action.}
Suppose now that a constituent is fixed. Choose
\(u\in\mathbf F_q^*\) with \(u^2=\lambda\), and put
\[
       b=u^{-1}a\in H,\qquad
       c=t^{(p^e-1)/2},\qquad
       T=c b\sigma_e\in GL_{\mathbf F_p}(A).
\tag{F1T37.13}
\]
Multiplication by the scalar \(u^{-1}\) changes no
conjugation on \(H\); neither does multiplication by
\(c\). Thus \(ThT^{-1}=\alpha(h)\).
The exact alternating form on the underlying
\(2rf\)-dimensional \(\mathbf F_p\)-space is
\(B_t(v,w)=\operatorname{Tr}_{q/p}(t\langle v,w\rangle)\).
It is nondegenerate, and (F1T37.13) preserves it, since
\[
 \operatorname{Tr}_{q/p}(t c^2 z^{p^e})
       =\operatorname{Tr}_{q/p}((tz)^{p^e})
       =\operatorname{Tr}_{q/p}(tz).
\tag{F1T37.14}
\]
Use (F1T37.6) to implement
\(D_t(v)\mapsto D_t(Tv)\) by a unitary \(U_T\).
As before, perfectness gives the exact intertwining
\(U_TR_t(h)U_T^{-1}=R_t(\alpha(h))\), and
\(U_T\Pi=\Pi U_T\). We have not assigned or changed
the scalar phase of any original inertia operator.
For a nonscalar original \(g\), (F1T37.3) is nonidentity,
so \(T\ne I,-I\).

\paragraph{6. The all-phase parity trace bound.}
Set \(N=2rf\),
\(a_+=\dim_{\mathbf F_p}\ker(T-I)\) and
\(a_-=\dim_{\mathbf F_p}\ker(T+I)\).
Conjugation by \(U_T\) permutes the exact orthogonal
basis (F1T37.6). Conjugation by \(\Pi U_T\) does so
through \(-T\). Taking their endomorphism traces gives
\[
       |\operatorname{Tr}U_T|^2=p^{a_+},\qquad
       |\operatorname{Tr}(\Pi U_T)|^2=p^{a_-}.
\tag{F1T37.15}
\]
As \(p\) is odd and \(T\ne\pm I\),
\(a_++a_-\le N\) and \(a_\pm\le N-1\).
The parity projector \((1+\varepsilon\Pi)/2\) therefore
gives
\[
 \bigl|\operatorname{Tr}_{\mathcal W_t^\varepsilon}U_T\bigr|
 \le \frac{p^{a_+/2}+p^{a_-/2}}{2}
 \le \frac{p^{(N-1)/2}+\sqrt p}{2}.
\tag{F1T37.16}
\]
For the last inequality, if one exponent is zero increase
it to one when possible; otherwise increase the sum of
exponents to \(N\). On that segment the convex function
\(p^{x/2}+p^{(N-x)/2}\), \(1\le x\le N-1\),
has its maximum at an endpoint. This also bounds the
case of an exponent zero directly by the same endpoint.

Write \(n=rf\ge5\). Dividing by the smaller dimension
in (F1T37.8) yields the uniform strict inequality
\[
 \frac{\bigl|\operatorname{Tr}_{\mathcal W_t^\varepsilon}U_T\bigr|}
       {d_\varepsilon}
 \le \frac{p^{n-1/2}+\sqrt p}{p^n-1}<\frac35.
\tag{F1T37.17}
\]
Here are rational checks valid for the full range. For
\(p=3\), the middle expression decreases as \(n\)
increases; \(1/\sqrt3<7/12\) and
\(\sqrt3<7/4\) give at \(n=5\) the upper bound
\(287/484<3/5\). For \(p\ge5\), use
\(1/\sqrt p<1/2\) and \(\sqrt p\le p/2\).
The resulting bound
\((p^n+p)/(2(p^n-1))\) is smaller than \(3/5\)
because \(p^n>5p+6\) for \(p\ge5,n\ge5\).
Multiplying any implementer by a unit scalar leaves
these inequalities unchanged.

\paragraph{7. All original multiplicities and inertia phases.}
Use the literal central idempotents and evaluation maps
of (F1T36.5)--(F1T36.6). Absolute irreducibility makes
the occurring irreducible \(H\)-modules a single orbit
under the original \(G\). If one constituent has level one,
(F1T37.12) makes all of them members of the four
modules above. In particular, all have the same parity
and the same degree \(d\); the possible exchange of
square-class labels is retained.

On a summand fixed by an original \(g\), the actual
evaluation isomorphism identifies its operator as
\[
  \rho(g)|_{\mathcal W_t^\varepsilon\otimes M_t}
        =U_T^\varepsilon\otimes B_{g,t},
 \qquad M_t=\operatorname{Hom}_H
               (\mathcal W_t^\varepsilon,V_{\mathbf C}).
\tag{F1T37.18}
\]
Indeed, after composing with \((U_T^\varepsilon)^{-1}\),
the operator commutes with the actual \(H\)-action;
Schur's lemma gives precisely the displayed endomorphism
of the Hom space. Give \(V_{\mathbf C}\) an original
\(G\)-invariant Hermitian metric. On an isotypic summand
it is the invariant metric on the irreducible factor
times a positive metric on \(M_t\), by the same
commutant calculation. The unitary choices above
therefore make \(B_{g,t}\) unitary. This operator
contains every original scalar phase; no phase is reset.
A summand moved to a different isotypic component
contributes zero to the original trace.

Using \(|\operatorname{Tr}B_{g,t}|\le\dim M_t\)
and (F1T37.17) on every fixed summand proves
\[
                  |\operatorname{Tr}_V(g)|<3m/5
             \qquad(g\notin Z).
\tag{F1T37.19}
\]
For a unitary original operator with
\(s=\operatorname{rank}(g-I)\), its fixed space
contributes \(m-s\) to the real trace and its other
\(s\) eigenvalues contribute at least \(-s\).
Thus \(s\ge(m-|\operatorname{Tr}g|)/2>m/5\)
for every nonscalar \(g\). A nonidentity scalar has
rank \(m\). Every original inertia is nonidentity,
so the original source budget gives
\[
 2m+6=\sum_{j=1}^{b}\operatorname{rank}(g_j-I)
           >\frac{b m}{5}\ge\frac{12m}{5},
\tag{F1T37.20}
\]
which forces \(m<15\), contrary to the retained
\(m\ge27\). This proves the asserted exclusion.

\paragraph{8. The original geometric factor and exact scope.}
All of (F1T37.18)--(F1T37.20) concern the original
representation of the original monodromy group on
the original curve cover. Inflation along \(\beta\),
the averaged central idempotents and the Hom evaluation
are the specified comparison maps; the group action
on their images is the original action conjugated
through those isomorphisms. The parabolic cochain
complex, its six-dimensional primitive corner and
its cup-product Hermitian pairing are unchanged by
the identity maps on that source and corner, with
polarization multiplier one. The same identity
preserves its Poincare correspondence, integral
lattice and projector, Hilbert/Chow cycle maps,
variational cycle-class maps and Adams--GRR maps
recorded in F01--F04. No tensor constituent is
substituted for the original coefficient and no new
detector component is counted. Formula (F1T37.1)
is a strict restriction of the existing T36 source
range; the other bounded levels remain unresolved
by this argument.

\paragraph{Sources and verification scope.}
The operative finite-group classification inputs are
Malle--Testerman, \emph{Linear Algebraic Groups and
Finite Groups of Lie Type} (2011), Theorem~24.17,
Remark~24.19 and Tables~24.2--24.3, and Steinberg,
\emph{Lectures on Chevalley Groups} (1968),
Theorem~30, finite-field proof steps~(1)--(4),
printed pp.~158--162. The exact level convention is
Guralnick--Larsen--Tiep, \emph{Inventiones mathematicae}
\textbf{235} (2024), Definition~3.2.
Guralnick--Magaard--Tiep,
\emph{Bulletin of the Institute of Mathematics,
Academia Sinica, New Series} \textbf{13} (2018),
443--461, pp.~444--445, identifies the published
Heisenberg convention and the two Weil labels.
The linear lift, irreducibility, parity, label
distinction, every semilinear intertwiner and the
uniform trace bound used here are proved above.
The separate source record distinguishes these
freshly read passages from their earlier cited
inputs; the supplemental finite certificate
checks explicit instances of the displayed maps.
